\scanpage{33}
% archive sheet number: 32
\inserted{(Déformation de schémas abéliens)}

\uline{Théorème 1} \quad $\deleted{\unsure{X|S}}S$ préschéma loc. noethérien \uline{connexe}

$X$ $S$-préschéma \uline{propre} et \uline{simple}

$e$ une section de $X$ sur $S$

$s\in S$ tel que $X_s=f^{-1}(s)$ soit un schéma abélien\deleted{\unsure{ne}} sur $k(s)$.

Alors il existe sur $X$ une seule structure de schéma abélien sur $S$, admettant $e$ comme section unité.

\uline{Cor.} \quad $\overline{\mathrm M}$ si $S$ n'est pas connexe, une structure abélienne est déterminée par la connaissance de sa section unité. (On va \unsure{le|la} \unsure{retrouver|utiliser})

\deleted{\uline{Théorème 2} \quad $S$ \deleted{\unsure{locam|local}} noeth.}

\deleted{$X_1$ et $X_2$ \inserted{de type fini} au-dessus de $S$ ($\varphi_i:X_i\to S$). On suppose}

\deleted{(i) les $\varphi_i$ plats et $\varphi_{2*}(\underline O_{X_2})=\underline O_S$, \uline{ou} $\varphi_2$ séparable}

\deleted{(ii) les fibres de $X_1$ connexes, ou $X_1$ est connexe.}

\deleted{$X_1$ muni d'une section $s_1$}

\deleted{\deleted{Alors \unsure{soit|tout}} $G$ schéma en groupes sur $S$, \deleted{abélien},}
\margin{?}

\deleted{$u:X_1\times_S X_2\longrightarrow G$, \inserted{\unsure{Considérons}} $u_2:X_2\longrightarrow G$,}

\deleted{$\exists\,u_1:X_1\longrightarrow S$ uniques, s.th.}
\[
\deleted{u=\deleted{\unsure{u_1+u_2+0}}u_1\mathrm{pr}_1+u_2\mathrm{pr}_2}
\]

\deleted{\uline{Cor.} \quad If $X_1=X_2$\deleted{\unsure{id. groupe}} \inserted{$=X$} \unsure{group scheme} over $S$, fibres are connected and simple, $X/S$ flat, $u:X\longrightarrow G$ s.th. $u(e)=e$, Then $u$ is a group hom.}

\diagram{\begin{tikzcd}[ampersand replacement=\&]
\deleted{X\times X}\arrow[r]\arrow[d,"\unsure{u{,}u|u\times u}"'] \& \deleted{X}\arrow[d,"u"]\\
G\times G\arrow[r] \& G
\end{tikzcd}}
